% PolyformTown — Extended Technical Summary
% Compile with: pdflatex -shell-escape summary.tex
%
\documentclass[10pt,a4paper]{article}

%% ---------- Packages ----------
\usepackage[T1]{fontenc}
\usepackage[utf8]{inputenc}
\usepackage{geometry}
\geometry{margin=2cm, top=2.2cm, bottom=2.5cm}

\usepackage{microtype}
\usepackage{parskip}
\usepackage{xcolor}
\usepackage{booktabs}
\usepackage{amsmath,amssymb}
\usepackage{graphicx}
\usepackage{hyperref}
\usepackage{enumitem}
\setlist{noitemsep, topsep=2pt}

%% minted — syntax highlighting (requires -shell-escape and pygments)
\usepackage{minted}
\setminted{
  fontsize=\footnotesize,
  baselinestretch=1.1,
  breaklines=true,
  linenos=false,
  frame=none,
  bgcolor=codebg
}
\definecolor{codebg}{HTML}{F5F5F5}

%% tcolorbox — side-by-side description/code panels
\usepackage[most]{tcolorbox}
\tcbuselibrary{minted,skins,breakable}

%% Colours
\definecolor{titlecolor}{HTML}{1A4D8F}
\definecolor{headercolor}{HTML}{2E6DB4}
\definecolor{rulecolor}{HTML}{AECBEF}

%% Side-by-side box style
\tcbset{
  dualbox/.style={
    enhanced,
    colframe=rulecolor,
    colback=white,
    fonttitle=\bfseries\small,
    breakable,
    boxrule=0.6pt,
    arc=3pt,
    sidebyside,
    sidebyside align=top seam,
    sidebyside gap=6mm,
    lefthand ratio=0.46,
    before skip=8pt,
    after skip=8pt,
  },
  descside/.style={
    colback=white,
  },
  codeside/.style={
    colback=codebg,
  }
}

%% Section colours
\usepackage{titlesec}
\titleformat{\section}{\large\bfseries\color{titlecolor}}{}{0em}{}[\vspace{-2pt}\color{rulecolor}\hrule\vspace{4pt}]
\titleformat{\subsection}{\normalsize\bfseries\color{headercolor}}{}{0em}{}

%% Convenient macros
\newcommand{\code}[1]{\texttt{\small #1}}
\newcommand{\file}[1]{\texttt{\footnotesize #1}}

%% ---------- Document ----------
\begin{document}

\begin{center}
  {\Huge\bfseries\color{titlecolor} PolyformTown}\\[4pt]
  {\large\color{headercolor} Extended Technical Summary — V5.0}\\[6pt]
  {\small A lattice polyform enumeration engine for hat-tiling research}\\[2pt]
  {\small\color{gray} \today}
\end{center}

\vspace{6pt}
\noindent\rule{\linewidth}{1pt}
\vspace{4pt}

% -----------------------------------------------------------------------
\section{Project Overview}

PolyformTown is a research-grade C engine that enumerates
\emph{lattice polyforms} — connected shapes built by joining tiles
along shared edges — and validates counts against OEIS reference
sequences.  The ultimate research goal is to prove properties of the
\textbf{hat tiling}, an aperiodic monotile discovered in 2023.

The engine supports three lattice types with a single shared geometric
core, and refines results through a layered \emph{run-level} (RL)
pipeline that combines combinatorial enumeration, canonical hashing,
and SAT-based verification.

\vspace{4pt}
\begin{tabular}{@{}lll@{}}
\toprule
\textbf{Lattice} & \textbf{Symmetry group} & \textbf{Example tiles} \\
\midrule
Square      & $D_4$ (8 transforms)  & monomino, domino, chair \\
Triangular  & $D_6$ (12 transforms) & triangle, hexagon, hh   \\
Tetrille    & $D_6$ tagged          & kite, hat               \\
\bottomrule
\end{tabular}

\vspace{6pt}
Some OEIS sequences the engine validates against:

\vspace{2pt}
\begin{tabular}{@{}lll@{}}
\toprule
\textbf{Sequence} & \textbf{Description} & \textbf{Lattice} \\
\midrule
A000104 & Polyominoes without holes & Square \\
A000105 & Free polyominoes          & Square \\
A000577 & Polyiamonds               & Triangular \\
A000228 & Polyhexes                 & Triangular \\
\bottomrule
\end{tabular}

% -----------------------------------------------------------------------
\section{Repository Layout}

\begin{tcolorbox}[dualbox, title=Directory Structure]
  The repository root contains a nested \file{PolyformTown/}
  project directory.  All source, data, tiles, and test artefacts live
  inside that subdirectory.

  \begin{description}[leftmargin=0pt, style=nextline]
    \item[\file{src/core/}] Geometry primitives: cycles, tiles,
      attachment, hashing, lattice transforms.
    \item[\file{src/throughput/}] Main growth pipelines:
      \code{poly\_pipeline}, \code{vcomp}.
    \item[\file{src/rl0/}--\file{src/rl6/}] Run-level modules; each
      refines the prior level's output.
    \item[\file{src/usage/}] CLI tools: \code{poly\_count},
      \code{poly\_print}, \code{vcomp\_count}, depict tools.
    \item[\file{tiles/}] Tile DSL files (\file{.tile}).
    \item[\file{data/}] Generated artefacts (not committed).
    \item[\file{tests/}] Smoke and regression scripts.
    \item[\file{manual/}] Documentation and theory notes.
  \end{description}
\tcblower
\begin{minted}{text}
PolyformTown/
+-- src/
|   +-- core/        # Coord, Cycle, Poly, Edge,
|   |                #   Tile, Hash, Lattice
|   +-- throughput/  # poly_pipeline, vcomp
|   +-- rl0/ .. rl6/ # run-level modules
|   +-- meta/        # dependency scanner
|   +-- usage/       # poly_count, poly_print ...
+-- tiles/
|   +-- hat.tile
|   +-- monomino.tile
|   +-- hexagon.tile  ...
+-- data/
|   +-- rl0/  rl1/  rl5/ ...
+-- tests/
|   +-- smoke.sh
|   +-- comprehensive.sh
+-- Makefile
+-- README.md
\end{minted}
\end{tcolorbox}

% -----------------------------------------------------------------------
\section{Core Data Structures}

\subsection{Coordinate and Cycle}

\begin{tcolorbox}[dualbox, title={Coordinate, Cycle, Poly, Edge --- \file{src/core/cycle.h}}]
  Every vertex in every lattice is represented by a \code{Coord}.
  The integer field \code{v} carries a \emph{valence tag}
  ($v \in \{6, 4, 3\}$) used exclusively in the tetrille lattice;
  for square and triangular lattices it is set to \code{4} or
  \code{6} respectively and acts as a system identifier.

  A \textbf{Cycle} is a closed polygon boundary stored as an ordered
  vertex array.  A \textbf{Poly} packages one outer cycle plus zero or
  more hole cycles.

  \begin{itemize}
    \item Outer boundary: \emph{counterclockwise} (CCW)
    \item Hole boundaries: \emph{clockwise} (CW)
    \item Maximum 384 vertices per cycle, 32 cycles per poly
  \end{itemize}

  The \textbf{signed area} (shoelace formula) determines orientation:
  positive $\Rightarrow$ CCW (outer), negative $\Rightarrow$ CW (hole).
\tcblower
\begin{minted}{c}
/* cycle.h */
typedef struct {
    int v;   /* valence tag: 6/4/3 */
    int x, y;
} Coord;

typedef struct {
    int n;              /* vertex count */
    Coord v[MAX_VERTS]; /* MAX_VERTS = 384 */
} Cycle;

typedef struct {
    int cycle_count;
    Cycle cycles[MAX_CYCLES]; /* [0]=outer, rest=holes */
} Poly;

typedef struct {
    Coord a, b;  /* directed edge a→b */
} Edge;

/* Shoelace area (×2) — sign gives orientation */
long long cycle_signed_area2(
    const Cycle *c, int lattice);
\end{minted}
\end{tcolorbox}

\subsection{Tile}

\begin{tcolorbox}[dualbox, title=Tile — \file{src/core/tile.h}]
  A \textbf{Tile} bundles the base boundary with all its pre-computed
  symmetry variants and the geometric metadata needed for SVG output.

  \begin{description}[leftmargin=0pt]
    \item[\code{base}] Canonical boundary loaded from the
      \file{.tile} file.
    \item[\code{variants}] Up to 12 symmetry images (rotations and
      reflections).  Computed once by \code{tile\_build\_variants()}.
    \item[\code{constants}] Named symbolic constants, e.g.\ $r =
      \sqrt{3}$, used in the basis matrix expressions.
    \item[\code{bases}] Per-valence $2{\times}2$ basis matrices that
      map integer lattice coordinates to real-plane coordinates for
      rendering.
  \end{description}

  Tile files use a simple DSL listing \code{cycle} vertices,
  optional \code{constants}, optional \code{basis} matrices,
  and the \code{lattice} type (\code{square}, \code{triangular}).
\tcblower
\begin{minted}{c}
/* tile.h */
#define TILE_LATTICE_SQUARE      0
#define TILE_LATTICE_TRIANGULAR  1
#define TILE_LATTICE_TETRILLE    2
#define MAX_VARIANTS 12

typedef struct {
    char name[64];
    Cycle base;
    int lattice;
    int variant_count;
    Cycle variants[MAX_VARIANTS];
    int constant_count;
    TileConstant constants[16];
    int basis_count;
    TileBasis bases[16];
} Tile;

/* hat.tile (excerpt) */
// lattice: tetrille
// constants: r = sqrt(3)
// basis 6:  r/2  -1/2
//           r/2   1/2
// cycle:
//   6 0 0   4 -1 0   3 -1 0 …
\end{minted}
\end{tcolorbox}

\subsection{Hash Table}

\begin{tcolorbox}[dualbox, title=Open-addressing Hash Table — \file{src/core/hash.h\,/\,hash.c}]
  After canonicalization each polyform is inserted into a
  \textbf{HashTable} to ensure each symmetry class is counted exactly
  once across the growth level.

  The hash function is FNV-1a over all cycle vertices (valence tag,
  $x$, $y$).  Collision resolution uses linear probing; the table
  doubles when load exceeds 70\,\%.

  \code{hash\_insert()} returns \code{1} if the shape was new
  (and stores it), or \code{0} if it was already present.
\tcblower
\begin{minted}{c}
/* hash.h */
typedef struct { Poly p; int used; } HashSlot;
typedef struct {
    HashSlot *slots;
    size_t size, count;
} HashTable;

/* hash.c — FNV-1a over vertex triples */
static unsigned poly_hash(const Poly *p) {
    unsigned h = 2166136261u;
    h ^= (unsigned)p->cycle_count;
    h *= 16777619u;
    for (int k = 0; k < p->cycle_count; k++) {
        const Cycle *c = &p->cycles[k];
        for (int i = 0; i < c->n; i++) {
            h ^= (unsigned)(c->v[i].v + 10000);
            h *= 16777619u;
            h ^= (unsigned)(c->v[i].x + 10000);
            h *= 16777619u;
            h ^= (unsigned)(c->v[i].y + 10000);
            h *= 16777619u;
        }
    }
    return h;
}

/* resize doubles table, re-hashes all entries */
/* insert returns 1=new, 0=duplicate */
int hash_insert(HashTable *h, const Poly *p);
\end{minted}
\end{tcolorbox}

% -----------------------------------------------------------------------
\section{Lattice Geometry}

\subsection{Three Lattice Types}

\begin{tcolorbox}[dualbox, title={Lattice Abstraction --- \file{src/core/lattice.h}, \file{tetrille.h}}]
  The three lattice types share a common interface so that the
  attachment and canonicalization algorithms are lattice-agnostic.
  Key per-lattice quantities:

  \begin{description}[leftmargin=0pt]
    \item[Transform count] Number of distinct symmetry group elements.
    \item[Direction count] Coordination number of each vertex.
    \item[Embed point] Maps integer lattice coordinates to the real
      plane (used only for SVG rendering and tie-breaking, never for
      enumeration).
  \end{description}

  The \textbf{tetrille} lattice uses \emph{tagged coordinates}.
  Each vertex carries a type tag $v \in \{6, 4, 3\}$
  corresponding to the three vertex valences of the kite-and-dart
  decomposition.  Translations must be lifted through the 6-system
  using exact integer formulas to avoid cross-tag arithmetic errors.
\tcblower
\begin{minted}{c}
/* lattice.h */
int lattice_transform_count(int lattice);
int lattice_direction_count(int lattice);
void lattice_embed_point(int lattice,
    Coord p, double *x, double *y);

/* tetrille.h — tagged translation in 6-system */
void tetrille_translate_cycle(
    Cycle *c, int m6, int n6);

/* cycle.c — tag-aware translation */
static void cycle_translate_tetrille(
    Cycle *c, int m6, int n6) {
  for (int i = 0; i < c->n; i++) {
    Coord *p = &c->v[i];
    if      (p->v == 6) { p->x += m6; p->y += n6; }
    else if (p->v == 4) { p->x += 2*m6; p->y += 2*n6; }
    else if (p->v == 3) {
        p->x += 2*m6 + n6;
        p->y += -m6  + n6;
    }
  }
}
\end{minted}
\end{tcolorbox}

\subsection{Signed Area and Orientation}

\begin{tcolorbox}[dualbox, title=Shoelace Formula — \file{src/core/cycle.c}]
  The signed area of a cycle is used to determine its orientation and
  to identify the outer boundary (maximum positive area) versus holes
  (negative area) after each tile attachment.

  For non-tetrille lattices the shoelace formula is applied directly
  to integer coordinates.  For the tetrille lattice each vertex is
  first embedded in a scaled real plane (scale factor 6) using fixed
  integer arithmetic to avoid floating-point errors.

  \[
    A_2 = \sum_{i=0}^{n-1} (x_i \, y_{i+1} - x_{i+1} \, y_i)
  \]
  $A_2 > 0$: CCW (outer boundary); \;
  $A_2 < 0$: CW (hole).
\tcblower
\begin{minted}{c}
/* cycle.c */
long long cycle_signed_area2(
    const Cycle *c, int lattice) {
  int tet = (lattice == TILE_LATTICE_TETRILLE);
  long long s = 0;
  for (int i = 0; i < c->n; i++) {
    Coord a = c->v[i];
    Coord b = c->v[(i + 1) % c->n];
    if (tet) {
      long long ax, ay, bx, by;
      tetrille_embed6_scaled(a, &ax, &ay);
      tetrille_embed6_scaled(b, &bx, &by);
      s += ax * by - bx * ay;
    } else {
      s += (long long)a.x * b.y
         - (long long)b.x * a.y;
    }
  }
  return s;
}
/* s > 0  →  CCW  (outer)   */
/* s < 0  →  CW   (hole)    */
\end{minted}
\end{tcolorbox}

% -----------------------------------------------------------------------
\section{Polyform Enumeration Pipeline}

\subsection{Growth Loop}

\begin{tcolorbox}[dualbox, title=Core Growth Loop — \file{src/throughput/poly\_pipeline.c}]
  The growth loop is the heart of the engine.  Starting from the
  canonical seed (a single tile), it iterates level by level:

  \begin{enumerate}
    \item Collect all boundary edges of every polyform in the current
      level.
    \item For each tile variant, try to attach it at each
      (frontier-edge, tile-edge) pair.
    \item Canonicalize the result under the lattice symmetry group.
    \item Filter with the live-boundary test (optional).
    \item Insert into a hash table; only new canonical shapes advance
      to the next level.
  \end{enumerate}

  A callback \code{on\_level()} fires at the start of each level and
  can terminate the loop early, making the pipeline reusable for both
  counting and printing tools.
\tcblower
\begin{minted}{c}
/* poly_pipeline.c */
void run_poly_levels(const Tile *tile,
    int max_n, int live_only,
    PolyLevelFn on_level, void *ud) {

  Poly seed_raw, seed;
  seed_raw.cycle_count = 1;
  seed_raw.cycles[0] = tile->base;
  poly_hash_key_lattice(&seed_raw,
      tile->lattice, &seed);

  PolyVec cur, next;
  vec_init(&cur, 64); vec_init(&next, 64);
  vec_push(&cur, &seed);

  for (int lv = 1; lv <= max_n; lv++) {
    if (on_level && !on_level(lv, &cur, tile, ud))
        break;
    if (lv == max_n) break;

    HashTable seen;
    hash_init(&seen, 1024);
    vec_clear(&next);

    for (size_t i = 0; i < cur.count; i++) {
      Edge fr[MAX_VERTS * MAX_CYCLES];
      int fc = build_boundary_edges(
                   &cur.data[i], fr);
      for (int v = 0; v < tile->variant_count; v++) {
        const Cycle *tv = &tile->variants[v];
        for (int be = 0; be < fc; be++)
         for (int te = 0; te < tv->n; te++) {
          Poly grown, canon;
          if (!try_attach_tile_poly(&cur.data[i],
               tv, tile->lattice, be, te, &grown))
              continue;
          poly_hash_key_lattice(&grown,
               tile->lattice, &canon);
          if (live_only &&
              !poly_has_live_boundary(&canon,tile))
              continue;
          if (hash_insert(&seen, &canon))
              vec_push(&next, &canon);
        }
      }
    }
    hash_destroy(&seen);
    PolyVec tmp = cur; cur = next; next = tmp;
  }
}
\end{minted}
\end{tcolorbox}

\subsection{Tile Attachment}

\begin{tcolorbox}[dualbox, title=Attachment Algorithm — \file{src/core/attach.c}]
  \code{try\_attach\_tile\_poly\_ex\_status()} is the geometric
  workhorse.  Given a frontier edge on the existing polyform and an
  edge of a tile variant, it:

  \begin{enumerate}
    \item \textbf{Aligns} the tile so the chosen tile edge is
      antiparallel and coincident with the chosen frontier edge.  For
      tetrille this requires lifting the translation into the 6-system.
    \item \textbf{Merges} all directed edges from both the polyform and
      the aligned tile into a flat list.
    \item \textbf{Cancels} antiparallel coincident edge pairs
      (the shared edge disappears from the boundary).  Parallel
      coincident edges indicate an illegal overlap and abort.
    \item \textbf{Walks} the remaining edges to extract boundary cycles
      via a left-first (outer) or right-first (hole) traversal.
    \item \textbf{Classifies} cycles by signed area and checks that no
      interior intersection occurs.
  \end{enumerate}

  Return codes (\code{AttachStatus}) distinguish geometry errors,
  buffer overflows, and valid results.
\tcblower
\begin{minted}{c}
/* attach.h */
typedef enum {
  ATTACH_STATUS_OK            =  1,
  ATTACH_STATUS_GEOMETRY      =  0, /* overlap */
  ATTACH_STATUS_BOUNDARY_BOUND= -1, /* too many edges */
  ATTACH_STATUS_CYCLE_BOUND   = -2, /* walk failed */
  ATTACH_STATUS_INTERNAL_BOUND= -3  /* interior hit */
} AttachStatus;

/* attach.c — edge cancellation core */
static int build_union_edges(
    const Poly *a, const Cycle *b,
    LEdge *out, int *out_n,
    AttachStatus *status) {
  /* collect a's edges, then b's edges */
  for (int i = 0; i < n; i++) {
    if (out[i].canceled) continue;
    for (int j = i+1; j < n; j++) {
      if (out[j].canceled) continue;
      if (edge_same(out[i].e, out[j].e)) {
        /* parallel duplicate → invalid */
        attach_status_set(status,
            ATTACH_STATUS_GEOMETRY);
        return 0;
      }
      if (edge_opp(out[i].e, out[j].e)) {
        /* antiparallel → cancel both */
        out[i].canceled = 1;
        out[j].canceled = 1;
        break;
      }
    }
  }
  return 1;
}
\end{minted}
\end{tcolorbox}

\subsection{Cycle Walking}

\begin{tcolorbox}[dualbox, title=Directed Cycle Extraction — \file{src/core/attach.c}]
  After edge cancellation the remaining directed edges form one or more
  closed cycles.  The walker picks the globally-lowest starting edge
  (by embedded position), then follows a \emph{prefer-left} path to
  trace the outer boundary and a \emph{prefer-right} path for holes.
  Direction scoring is modular-arithmetic over the lattice direction
  count, so it handles all three lattice types uniformly.
\tcblower
\begin{minted}{c}
/* attach.c — direction-scored next-edge picker */
static int pick_next_edge(
    const Edge *edges, int m,
    const int *used, Coord v,
    int prev_dir, int prefer_left,
    int lattice) {
  int dc  = lattice_direction_count(lattice);
  int rev = (prev_dir + dc/2) % dc;
  int best = -1;
  int best_score = prefer_left ? -1 : dc+1;

  for (int i = 0; i < m; i++) {
    if (used[i]) continue;
    if (!coord_eq(edges[i].a, v)) continue;
    int d = lattice_direction_index(
                lattice, edges[i]);
    if (d < 0) continue;
    int score = (d - rev + dc) % dc;
    if (score == 0) continue;
    if (prefer_left ? (score > best_score)
                    : (score < best_score)) {
      best_score = score; best = i;
    }
  }
  return best; /* -1 if no valid edge found */
}
\end{minted}
\end{tcolorbox}

% -----------------------------------------------------------------------
\section{Canonicalization}

\begin{tcolorbox}[dualbox, title=Canonical Form — \file{src/core/cycle.c\,/\,cycle.h}]
  Two polyforms are \emph{equivalent} if one can be obtained from the
  other by a lattice symmetry (rotation and/or reflection).  To count
  equivalence classes rather than labelled shapes, every new polyform
  is reduced to its \textbf{canonical representative}:

  \begin{enumerate}
    \item Apply all symmetry transforms ($|D_4|{=}8$ or
      $|D_6|{=}12$ elements).
    \item For each image, normalize position (translate so the
      lexicographically smallest vertex is at the origin).
    \item Rotate each cycle to its lexicographically minimal starting
      vertex.
    \item Sort hole cycles lexicographically.
    \item Select the image that is lexicographically smallest
      overall.
  \end{enumerate}

  \code{poly\_hash\_key\_lattice()} wraps this into the one call
  made by the growth loop before every hash insert.
\tcblower
\begin{minted}{c}
/* cycle.h */
void poly_hash_key_lattice(
    const Poly *src, int lattice, Poly *key);

/* Internally calls poly_canonicalize_lattice(),
   which applies all lattice transforms and
   keeps the lexicographic minimum. */

void poly_canonicalize_lattice(
    const Poly *src, Poly *out, int lattice) {
  int nt = lattice_transform_count(lattice);
  Poly best;
  int have_best = 0;
  for (int t = 0; t < nt; t++) {
    Poly tmp;
    poly_transform_lattice(src, &tmp,
        lattice, t);
    poly_normalize_position(&tmp, lattice);
    /* shift each cycle to its min vertex */
    for (int k = 0; k < tmp.cycle_count; k++)
        cycle_canonicalize_shift(
            &tmp.cycles[k]);
    /* sort hole cycles */
    /* compare to running best */
    if (!have_best || poly_less(&tmp, &best)) {
        best = tmp; have_best = 1;
    }
  }
  *out = best;
}
\end{minted}
\end{tcolorbox}

% -----------------------------------------------------------------------
\section{Vertex Completion (VComp)}

\begin{tcolorbox}[dualbox, title=Vertex-Completion Enumeration — \file{src/throughput/vcomp.h}]
  An alternative enumeration mode tracks how many distinct ways a
  specific \textbf{target vertex} can be completely surrounded by
  tiles.  This corresponds to enumerating valid \emph{vertex figures}
  and is critical for the run-level analysis.

  The state records:
  \begin{description}[leftmargin=0pt]
    \item[\code{target}] The vertex being completed.
    \item[\code{hidden}] Vertices discovered in the interior of the
      growing figure (constrained by the target's local geometry).
    \item[\code{ports}] Exposed attachment sites on the boundary that
      can still receive tiles.
    \item[\code{tiles}] History of attached tiles (when tracking is
      enabled).
  \end{description}

  Multiple levels record partial completions graded by the number of
  hidden vertices — a slowly-growing quantity that makes depth-limited
  search tractable.
\tcblower
\begin{minted}{c}
/* vcomp.h */
typedef struct {
  Poly poly;       /* current boundary shape */
  Coord target;    /* vertex being completed */
  Coord hidden[VCOMP_MAX_HIDDEN];
  int   hidden_count;
  Coord ports[VCOMP_MAX_HIDDEN]; /* open sites */
  int   port_count;
  Cycle tiles[VCOMP_MAX_TRACE];  /* history */
  int   tile_count;
} VCompRawState;

/* Enumerate all partial/complete vertex
   figures up to max_hidden interior verts */
void vcomp_enumerate_levels(
    const Poly *base,
    const Tile *tile,
    Coord target,
    const Coord *initial_hidden,
    int initial_hidden_count,
    const Coord *initial_ports,
    int initial_port_count,
    const Cycle *initial_tiles,
    int initial_tile_count,
    int max_hidden,
    int track_tiles,
    VCompLevels *out);
\end{minted}
\end{tcolorbox}

% -----------------------------------------------------------------------
\section{Run-Level Pipeline}

The run-level system progressively narrows the space of hat-tiling
configurations by applying successively stronger local constraints.
Each level produces a \emph{completions} dataset and optionally a
\emph{deletions} dataset consumed by the next level.

\subsection{RL0: Vertex Figures}

\begin{tcolorbox}[dualbox, title=RL0 — Vertex Figure Generation — \file{src/rl0/}]
  RL0 exhaustively enumerates all valid ways to tile completely around
  a single vertex.  Results are stored in \file{data/rl0/completions.dat}
  as multi-line records delimited by \code{---[N]---}.

  The \textbf{ForgetMap} (remembrance dictionary) indexes every
  vertex figure by its \emph{arc} — a cyclic sequence of
  (parity, tile-index) pairs — enabling rapid lookup of valid
  completions given any partial arc.  This dictionary is the primary
  filter used by RL1.

  The \textbf{deletion set} records vertex figures that are rejected
  by refinement passes (e.g.\ because they cannot be extended to a
  valid surround at the next level).
\tcblower
\begin{minted}{c}
/* forget_map.h */
typedef struct { int p; int i; } RL0FMItem;
/* p = parity (±1), i = tile variant index */

typedef struct {
  int n;
  RL0FMItem item[16]; /* arc of (p,i) pairs */
} RL0FMCycle;         /* also used as RL0FMArc */

typedef struct {
  RL0FMArc key;
  RL0FMArc *values;
  int value_count;
  int value_capacity;
} RL0FMRow;

typedef struct {
  RL0FMCycle *cycles;
  int cycle_count, cycle_capacity;
  RL0FMRow   *rows;
  int row_count, row_capacity;
} RL0ForgetMap;

/* Look up all valid completions for a partial arc */
int rl0_fm_lookup(const RL0ForgetMap *map,
    const RL0FMArc *key,
    const RL0FMArc **values,
    int *value_count);
\end{minted}
\end{tcolorbox}

\subsection{RL1–RL3: Surround Levels}

\begin{tcolorbox}[dualbox, title=RL1–RL3 — Boundary Completion Filters — \file{src/rl1/} \ldots \file{src/rl3/}]
  Each surround level takes the completion records from the previous
  level and checks whether each vertex figure can be extended to a
  valid single (RL1), double (RL2), or triple (RL3) ring of surrounding
  tiles.

  Extension uses a DFS-based \textbf{boundary completion} (\code{bcomp})
  search that respects the RL0 remembrance dictionary: every partial
  vertex figure encountered during the extension must appear in the
  dictionary, or the branch is pruned immediately.

  Records that \emph{cannot} be extended are added to the
  deletions file, which is consumed by the next level's loader
  (via \code{--depth} refinement).

\tcblower
\begin{minted}{bash}
# Build pipeline (Makefile)
make boot          # RL0 → RL5 in sequence
make boot RL3_REFINE_DEPTH=2  # deeper pruning

# RL0
./bin/rl0_generate [tilefile]
./bin/rl0_refine

# RL1
./bin/rl1_generate \
    --rl0 data/rl0/completions.dat \
    --output data/rl1/completions.dat

./bin/rl1_refine \
    --input  data/rl1/completions.dat \
    --output data/rl1/deletions.dat \
    --depth 1

# RL3 with deeper elimination
./bin/rl1_refine \
    --input  data/rl3/completions.dat \
    --output data/rl3/deletions.dat \
    --depth 2
\end{minted}
\end{tcolorbox}

\subsection{RL4: Supertile Extraction}

\begin{tcolorbox}[dualbox, title=RL4 — Supertile Candidates — \file{src/rl4/}]
  RL4 extracts candidate \emph{supertile} shapes from the RL3-filtered
  surround set.  A supertile is a second-order metatile whose internal
  structure is fully determined by the local hat-tiling constraints
  established in RL0–RL3.

  Two supertile preferences are written:
  \begin{description}[leftmargin=0pt]
    \item[\file{preferences/focus.supertile}] Smaller default
      candidate used in subsequent RL5 analysis.
    \item[\file{preferences/overlap.supertile}] Maximal overlap
      candidate for cross-checking.
  \end{description}
\tcblower
\begin{minted}{bash}
# Regenerate RL4 data
make rl4_refine_data
make rl4_refine_svg

# Override supertile for RL5 extraction
make rl4_supertile_hexagons_data \
  RL4_SUPERTILE=preferences/overlap.supertile \
  RL4_SUPERTILE_HEXAGONS_DATA=\
      data/rl4/overlap_supertile_hexagons.dat
\end{minted}
\end{tcolorbox}

\subsection{RL5: Hexagon Model and SAT Verification}

\begin{tcolorbox}[dualbox, title=RL5 — Hexagon Model — \file{src/rl5/hex5.h}]
  RL5 operates on an abstract \textbf{hexagon model}: each hat
  supertile is treated as a labelled hexagon, with edge labels
  encoding which orientations of the hat tile can be placed on each
  face.  Adjacency constraints become edge-matching rules.

  The \code{Hex5Model} holds:
  \begin{itemize}
    \item \code{tiles} — all canonical hexagon label sets.
    \item \code{oriented} — all symmetry images.
    \item \code{rule\_a / rule\_b} — edge-compatibility rules
      (which label on face $a$ can be adjacent to label on face $b$).
  \end{itemize}

  \code{BComp5Result} stores complete boundary-completion rings
  found by the DFS over the hexagon constraint graph.  The
  \code{Inner5Set} records pairs of hexagon labels that appear
  back-to-back in any completed ring — used as inner boundary
  constraints.
\tcblower
\begin{minted}{c}
/* hex5.h */
#define HEX5_SIDES 6
typedef struct {
    char e[HEX5_SIDES][64]; /* 6 edge labels */
} Hex5;

typedef struct {
    Hex5 tiles[512];     int ntiles;
    Hex5 oriented[4096]; int noriented;
    char rule_a[4096][64];
    char rule_b[4096][64];
    int nrules;
} Hex5Model;

typedef struct {
    int  center;
    int  r[HEX5_SIDES]; /* ring of 6 tiles */
} Ring5;

typedef struct {
    Ring5 rings[262144]; int nrings;
    long starts, closure_tests, closure_hits;
} BComp5Result;

/* Build vertex-completion lookup */
int vcomp5_build(
    const Attach5Dict *a, VComp5Dict *v);

/* Build boundary-completion rings */
int bcomp5_build(const Hex5Model *m,
    const VComp5Dict *v, BComp5Result *b);
\end{minted}
\end{tcolorbox}

\begin{tcolorbox}[dualbox, title=RL5 — SAT Verification — \file{src/rl5/sat5.c}]
  For pathological supertile configurations that cannot be eliminated
  by local constraint propagation alone, RL5 generates a SAT instance
  and calls an external solver (pysat / PySAT library) to determine
  satisfiability.

  The SAT encoding maps each possible tile placement at each position
  to a Boolean variable; the local adjacency and vertex-completion
  constraints become clauses in CNF form.  A SAT result of
  \emph{unsatisfiable} proves that the configuration cannot appear in
  any valid hat tiling.

  Output artefacts for both supertile candidates:
  \begin{itemize}
    \item \file{data/rl5/hexagons.dat}
    \item \file{data/rl5/supertile\_hexagons.dat}
    \item \file{data/rl5/overlap\_supertile\_hexagons.dat}
    \item \file{data/rl5/equivalence\_candidates.dat}
  \end{itemize}
\tcblower
\begin{minted}{bash}
# Generate RL5 hexagon datasets
make rl5_generate

# Run equivalence scaffold
make rl5_refine_data

# Inspect candidates
less data/rl5/equivalence_candidates.dat

# Pipe SAT instance to solver (example)
./bin/rl5_refine \
    --hex  data/rl5/hexagons.dat \
    --out  data/rl5/equivalence_candidates.dat

# SAT invocation (internal)
# python3 -m pysat.examples.rc2 <cnf>
\end{minted}
\end{tcolorbox}

% -----------------------------------------------------------------------
\section{Tile File Format}

\begin{tcolorbox}[dualbox, title=Tile DSL — \file{tiles/*.tile}]
  Tile files are plain-text DSL files specifying the tile name,
  lattice type, optional symbolic constants ($\sqrt{3}$, etc.),
  optional rendering basis matrices, and the boundary cycle as a list
  of tagged integer coordinates.

  The three example tiles below span all three lattice types.
  The hat tile uses all three coordinate valences in a single cycle,
  illustrating the complexity of the tetrille embedding.
\tcblower
\begin{minted}{text}
# monomino.tile (square lattice)
name: monomino
lattice: square
basis:
4:
  1 0
  0 1
cycle:
4 0 0   4 1 0   4 1 1   4 0 1

# triangle.tile (triangular lattice)
name: triangle
lattice: triangular
constants:
r = sqrt(3)
basis:
6:
  1 0
  1/2 r/2
cycle:
6 0 0   6 1 0   6 0 1

# hat.tile (tetrille — three valence systems)
name: hat
lattice: tetrille
constants:
r = sqrt(3)
cycle:
6  0  0   4 -1  0   3 -1  0
4 -1 -1   6  0 -1   4  0 -3
3 -1 -2   4  1 -3   3  0 -2
4  1 -2   6  1 -1   4  2 -1
3  1 -1   4  1 -1
\end{minted}
\end{tcolorbox}

% -----------------------------------------------------------------------
\section{CLI Tools and Output}

\begin{tcolorbox}[dualbox, title=Usage Examples]
  The build produces several CLI tools in \file{bin/}.  All tools
  accept an optional tile file argument and respect a \code{--live-only}
  flag that restricts output to polyforms whose boundary can still
  accept further tiles.

  Output can be piped into \code{imgtable} to produce SVG grids,
  enabling visual inspection of enumerated shapes at any level.

  Selected OEIS validations:
  \vspace{4pt}
  \begin{tabular}{@{}lll@{}}
  \toprule
  \textbf{Count} & \textbf{n} & \textbf{Target} \\
  \midrule
  Monomino & 7 & 107 \\
  Domino   & 5 & 4 \\
  Triangle & 8 & 66 \\
  Hexagon  & 6 & 82 \\
  \bottomrule
  \end{tabular}
\tcblower
\begin{minted}{bash}
# Count polyforms (square lattice)
./bin/poly_count 7 tiles/monomino.tile
./bin/poly_count 5 tiles/domino.tile

# Count polyforms (triangular lattice)
./bin/poly_count 8 tiles/triangle.tile
./bin/poly_count 6 tiles/hexagon.tile

# Vertex-completion count (hat / tetrille)
./bin/vcomp_count 3 tiles/hat.tile

# Pipe to SVG renderer
./bin/vcomp_print 3 tiles/hat.tile \
    | ./bin/imgtable > out.svg

# RL0 depiction grid
./bin/rl0_depict \
    --data data/rl0/completions.dat \
    --valence 6 --tile-count 3 --grouped \
    | ./bin/imgtable > img/rl0_v6_n3.svg

# RL0 refinement graph
./bin/rl0_refine --print \
    | ./bin/depict_search0_graph \
        --colors > img/rl0_refine.svg
\end{minted}
\end{tcolorbox}

% -----------------------------------------------------------------------
\section{Testing}

\begin{tcolorbox}[dualbox, title=Test Suite — \file{tests/}]
  The project includes two test scripts that validate core
  functionality and count parity between the two enumeration modes
  (edge-attachment and vertex-completion).

  \textbf{Smoke test} checks that:
  \begin{itemize}
    \item All CLI tools build and run without error.
    \item Basic polyomino counts match known OEIS values.
    \item Canonicalization is idempotent.
  \end{itemize}

  \textbf{Comprehensive test} additionally verifies:
  \begin{itemize}
    \item \code{vcomp\_parity}: tracked and untracked vcomp
      enumeration yield identical counts for domino, chair, and hat.
    \item Hole-count checks via \code{poly\_print} grep.
    \item RL0 remembrance accounting consistency.
  \end{itemize}
\tcblower
\begin{minted}{bash}
# Smoke test (fast)
bash tests/smoke.sh

# Full regression suite
bash tests/comprehensive.sh

# Expected hole counts from poly_print
# Monomino at level 7
./bin/poly_print 7 tiles/monomino.tile \
    | grep '^1 ' | wc -l   # → 107

# Domino at level 4
./bin/poly_print 4 tiles/domino.tile \
    | grep '^1 ' | wc -l   # → 3

# RL0 remembrance count check
make check_rl0_remembrance_counts
# writes data/rl0/remembrance_counts.txt
\end{minted}
\end{tcolorbox}

% -----------------------------------------------------------------------
\section{Mathematical Background}

\subsection{Polyform Counting}

A \emph{polyform} of size $n$ on a lattice is an equivalence class of
connected unions of $n$ tiles under the lattice symmetry group.  The
central algorithmic challenge is that the number of distinct polyforms
grows exponentially in $n$, with a growth rate (connective constant)
that depends on the tile geometry.

For polyominoes the growth constant is approximately $\lambda \approx
4.0626$, giving:
\[
  |\mathcal{P}_n| \sim C \cdot \lambda^n, \qquad \lambda \approx 4.0626.
\]

\subsection{Edge Cancellation and Euler Characteristic}

The boundary representation is closed under the edge-cancellation
operation.  After attaching a tile, the shared edge is cancelled; the
remaining directed edges form a valid set of boundary cycles whose
Euler characteristic satisfies:
\[
  V - E + F = 2 - 2g - h
\]
where $g$ is the genus and $h$ the number of holes.  In practice all
polyforms have genus 0, so the formula reduces to $F = 2 - V + E + h$.

\subsection{The Hat Tile and Aperiodicity}

The \textbf{hat} is a 13-sided polygon on the tetrille (kite-and-dart)
lattice whose discovery in 2023 provided the first known aperiodic
monotile: a single shape that tiles the plane but only aperiodically.

PolyformTown's run-level pipeline targets a computational proof of
key local properties of the hat tiling, specifically that certain
vertex figure constellations are forbidden, which would imply the
aperiodicity constraint.

% -----------------------------------------------------------------------
\section{Architecture Summary}

\begin{center}
\begin{tcolorbox}[
  colframe=rulecolor, colback=white, boxrule=0.6pt, arc=3pt,
  width=0.96\linewidth
]
\begin{minted}{text}
  +------------------------------------------+
  |          Tile DSL  (.tile)               |
  +------------------------------------------+
                 | tile_load() / tile_build_variants()
  +------------------------------------------+
  |        Core Layer  (src/core/)           |
  |  Coord  Cycle  Poly  Edge  Hash  Lattice |
  +--------------------+---------------------+
                       |
          +------------+------------+
          |                         |
  +-------+--------+    +-----------+---------+
  | poly_pipeline  |    |  vcomp_pipeline     |
  | (edge growth   |    |  (vertex-figure     |
  |  + hash)       |    |   completion)       |
  +-------+--------+    +-----------+---------+
          |                         |
  +-------+-------------------------+---------+
  |           Run-Level Pipeline              |
  |  RL0 -> RL1 -> RL2 -> RL3 -> RL4 -> RL5 |
  |  (vertex  (surrounds)  (supertile)        |
  |  figures)              (SAT verify)       |
  +-------------------------------------------+
\end{minted}
\end{tcolorbox}
\end{center}

\vspace{4pt}
\noindent\rule{\linewidth}{0.6pt}
\vspace{2pt}

{\small\color{gray}
PolyformTown V5.0 — stable, reproducible state with verified square,
triangular, and tetrille workflows.  SAT-based verification (RL5)
and equivalence scaffold are current research frontiers.
Compile with \texttt{pdflatex -shell-escape summary.tex} (requires
\texttt{pygments} for \texttt{minted}).
}

\end{document}
